\section{问题分析}

\subsection{数据预处理}
普通空白值在未经核验前视为不可裁决缺失，不自动替换为零。仅当互斥进流与出流字段满足“本行孤立事件、另一列明确非零，且备注、相邻记录和全表口径共同支持互斥解释”时，才将空白认定为结构零；所有处理均保留字段来源与审计记录。
本题同时包含任务级离散决策、分钟级连续占用和小时级能源结算。若把任务直接压缩为小时平均负荷，会漏掉只在小时内短暂重叠却造成瞬时 GPU 峰值的任务；若把所有触及某小时的任务都按整小时计费，又会高估 AI 电量。因此本文先保留任务的真实起止区间，再用“是否正重叠”审计瞬时容量，用“实际重叠时长”核算 GPU-hour 和电量，最后把重建负荷传给能源模型。

\subsection{问题一：预测与末端调度分离}
预测对象是区域—任务类型—小时的到达 GPU 需求。训练、验证、测试必须按时间顺序切分，且多步预测时后续滞后项只能使用此前预测值，不能回填测试真值。日周期与周周期朴素模型提供可解释基线，Huber 回归处理尖峰需求并通过验证集选择超参数\cite{huber1964robust}。预测只用于精度评价；最后 24 小时调度仍读取第 2376 至 2399 小时实际到达的 538 个任务，从而避免“以预测量替代真实任务”的口径错误。

末端任务可能跨越第 2399 小时，因此调度时域延伸到第 2405 小时；第 2406 小时只用于能源和 SOC 结算。实时推理必须到达即开工，弹性任务可在到达与截止之间选择整数开工时刻。候选区域还需同时通过单向时延、GPU、IT 功率和设施功率下界筛选。

\subsection{问题二：同一 CandidateID 流上的公平比较}
问题二的首要风险是把“任务触及某小时的瞬时功率”误作该小时整小时电量。本文分别维护瞬时 GPU、瞬时 IT 功率、按重叠时长累计的 AI 电量和由电量换算的设施小时平均功率。五个直接方向从同一 Q1 基线出发，并按目标无关的 round-robin 分层顺序读取完全相同的 CandidateID 前缀；前 50000 行先覆盖全部任务，随后两层再扩展同一任务的时间与区域候选。直接碳方向与额外多起点强化分开，后者不进入公平横向表。

\subsection{问题三：正式平衡负荷上的储能比较}
问题三读取五个 Q2 direct 方案的正式哈希，并按预先声明的下游合同把 formal balanced 的设施负荷作为共同输入。储能候选仍包含 no-action、成本、碳、新能源、峰值、爬坡和平衡七种策略；轨迹先去重和候选内支配筛选，再以统一加权理想点距离选择，避免按方案名称指定答案。

\subsection{问题四：按情景动态重建邻域}
联合问题中，任务方案改变设施负荷，储能轨迹又改变任务移动的边际电价和边际碳信号。固定若干预制任务方案再逐情景挑选，会切断这种反馈。本文对每个情景从唯一共同任务基线和真实初始 SOC 独立重启；不存在预制的两个替代任务方案。对当前负荷分别计算成本、平衡、新能源三类储能政策，按 Q4 自身的加权理想点规则选择，再依据所选轨迹形成价格、碳、购电和弃电信号，动态生成 600 个任务候选。接受改进后重建下一轮邻域。最多 8 轮，连续 2 轮无改善即停止。

17 个情景覆盖基准、4 个碳约束水平、3 个峰谷价差、3 个售电机制、3 个新能源水平，以及种子 2026、2027、2028 的新能源波动。正常联合基准若为正碳，四档碳目标直接以该基准定义；若为零碳，才改用新能源缩放 0.8 的短缺压力基准，且不与正常基准混淆。每类情景输出唯一任务调度数、唯一储能轨迹数、机制活跃性和退化原因。碳目标未达到只能表述为“声明启发式内未达到”，不能据此推断数学不可行。

\subsection{四问衔接与证据边界}
四问共享任务重叠、设施负荷、能源守恒和审计函数，但不直接拼接决策结果。问题二以问题一可行基线启动；问题三与问题四均读取修复后的五个 direct 正式哈希，并以问题二 formal balanced 的设施负荷为共同基线，统一使用 21 状态网格，但在每个情景和每轮独立重选三类储能政策。输出可信性来自结构化结果、逐行 CSV、完整约束审计和两次隔离复算。由于不存在独立 Oracle，最高可信等级限定为 scenario-feasible。
